\scanpage{19}
Coniques circonscrites à un polygone \unsure{regulier|régulier}.

Faisons les calculs en coordonnées \unsure{universelle|universelles}

\[
\left\{\begin{aligned}
s_0&=(0,0)\\
s_1&=(1,0)\\
s_{-1}&=(0,1)\\
s_2&=(1+\alpha,1)\\
s_{-2}&=(1,1+\alpha)
\end{aligned}\right. \tag*{(1)}
\]

Écrivons une fonction quadratique \unsure{général|générale} en coord. affines
\[
q(x,y)=ax^2+bxy+cy^2+ux+vy+d
\]
\deleted{\unsure{pour}} Écrivons que \deleted{\unsure{le lieu}} la conique qu'elle définit
passe par les \inserted{5 pts} $s_i$, $-2\le i\le +2$ ; on trouve

\[
\left.\begin{aligned}
q(s_0)&=d=0\\
q(s_1)&=a+u=0&\text{i.e. }u&=-a\\
q(s_{-1})&=\deleted{\unsure{b}}\unsure{c|b}+v=0&\text{--- }\ v&=-b
\end{aligned}\right\}
q(x,y)=a(x^2-x)+c(y^2-y)+bxy\quad\deleted{\unsure{+c}}
\]

\[
\left\{\begin{aligned}
q(s_2)&=a(1+\alpha)\alpha+b(1+\unsure{2|\alpha})=(1+\alpha)(\alpha a+b)=0\\
q(s_{\unsure{1|-2}})&=\deleted{\unsure{b}}c(1+\alpha)\alpha+b(1+\alpha)=(1+\alpha)(\alpha\deleted{\unsure{b}}c+b)=0
\end{aligned}\right.
\]

Si $\alpha-1$ inv. (i.e. on n'a de triangle sur aucune fibre
la condition s'écrit
\[
 b=-\alpha a\inserted{\unsure{=}}\deleted{\underbrace{\unsure{,\,a=-\alpha b}}}\,=-\alpha c
\]
ou encore

et si $\alpha$ inversible (i.e. on n'a pas de carrés sur aucune fibre)
\[
 a=c,\quad b=-\alpha a,\quad \unsure{''}
\]
donc $q$ doit être multiple \inserted{$a q_0$} de la forme
\[
q_0(x,y)=(x^2-x)+(y^2-y)-\alpha xy \tag*{(2)}
\]
% slanted
\margin{\uline{NB} C'est la normalisation pour laquelle $q(s-s')=1$ si $s,s'$ sont 2 sommets consécutifs \unsure{de $P$|$(P)$} \ldots}

Comme celle-ci est invariante sous $G$, la

% archive sheet number: 18

